\section{Whole fibres and quantitative global distances}\label{sec:global}

The local estimates have relative errors, whereas gluing distances
introduces additive errors. We keep these separate. First we obtain
whole-fibre charts and the central completion. We then compare short
ambient pairs in both directions, and finally concatenate geodesics.

\subsection{Coverage and compatibility}
\begin{proposition}\label{prop:coverage}
Under \textup{(G1)--(G3)}, local classifying charts can be chosen to
contain the whole preimages of smaller target neighborhoods, for
all sufficiently small generic parameters. For each prescribed
relative accuracy they give finite covers with that accuracy.
Their coarse transitions satisfy all cocycle identities, and if
two squared-metric errors are \(\epsilon_i,\epsilon_j<1/2\), then
\begin{equation}\label{eq:overlap-metric}
 \frac{1-\epsilon_i}{1+\epsilon_j}g_i
 \le h_{ij,\theta}^*g_j
 \le\frac{1+\epsilon_i}{1-\epsilon_j}g_i .
\end{equation}
\end{proposition}

\begin{proof}
By (G2), a sequence of nearby chamber connections has a subsequence
converging smoothly modulo gauge to a central polystable bundle.
Indeed surface compactness first gives an ideal limit; the
small-energy estimates and removability leave only positive
integral lost charge, which (G2) excludes. The smoothly varying
background metrics and determinant connections preserve these
local estimates; see also \cite[Lemma~3.2]{wuebben17}.

Suppose \(\theta_n\to0\) and \(\pi_{\theta_n}(E_n)\to p\).
Let \(E\) be a smooth gauge subsequential limit.
Inverse coverage in the local solution chart at \(E\) gives
balanced coefficients \(x_n\to0\), after compact-group adjustment.
Equation~\eqref{eq:classifying-diagram} identifies their target
classes with the affine invariants of \(x_n\), which converge
to the vertex at \([E]\). Since \(Z\) is Hausdorff, \([E]=p\).

Fix an outer chart at \(p\). If no smaller target neighborhood
and parameter ball had their whole preimages inside it, choose
points outside the chart with targets tending to \(p\) and
parameters tending to zero. The preceding conclusion puts them
in its open actual solution neighborhood, a contradiction.
Finite invariant properness bounds balanced representatives with
bounded invariants and small levels. Shrinking the target neighborhood
therefore puts its whole preimage in any chosen inner coefficient
ball. Conversely, the finite projective quotient is onto its affine
target, and its small representatives classify actual nearby
bundles. This gives coverage on both sides.

For a prescribed accuracy, shrink an outer chart separately at
each \(p\) using \eqref{eq:relative}. Extract a finite subcover of
their inner target images, and only then shrink the common parameter
ball. This order avoids any claim of a uniform complementary
elliptic gap for separate stable backgrounds approaching a
reducible limit. A chart at the reducible limit covers that region.

On overlaps put \(h_{ij,\theta}=f_{j,\theta}^{-1}f_{i,\theta}\).
These are the unique coarse identifications of the same holomorphic
bundle classes, so they satisfy \(h_{jk,\theta}h_{ij,\theta}
=h_{ik,\theta}\). They are holomorphic in complex coordinates and
smooth on smooth quotient pieces. Comparing each finite metric
with the same actual harmonic metric gives
\eqref{eq:overlap-metric}.
\end{proof}

Universal families on overlapping domains can differ by a base
line or a gerbe presentation. The coarse transition used above is
unaffected. Proposition~\ref{prop:coverage} also supplies the
geometric cover used in Proposition~\ref{prop:descent}.

\subsection{The central length space}
Give each central representation-type stratum its harmonic metric.
Let \(d_0\) be the infimum of the lengths of paths with finitely
many smooth stratum pieces and finite one-sided endpoint lengths.

\begin{proposition}\label{prop:central-completion}
The function \(d_0\) is a metric inducing the projective topology
of \(Z\). It makes \(Z\) a compact geodesic space and is the
completion metric of the intrinsic harmonic metric on \(U\).
A finite upper bound \(D_*\ge\diam(Z,d_0)\) can be constructed
from one finite central chart cover and its coarse metric bounds.
\end{proposition}

\begin{proof}
At zero level the model radius is the distance to the cone vertex.
Small model neighborhoods thus connect nearby points at cost
bounded by a fixed multiple of their radii. The local tensor
comparison transfers these paths to \(Z\). Connectedness and a
finite chart cover make all distances finite.

There is also a lower bound for leaving a chart. With squared
metric comparison between \(1/2\) and \(2\), a path beginning at
model radius \(r<R\) and leaving the radius-\(R\) sublevel has
length at least
\begin{equation}\label{eq:exit}
 (R-r)/\sqrt2 .
\end{equation}
The model radius is 1-Lipschitz, and the image of a closed sublevel
is compact, so a first exit must meet its radial boundary.
This applies to ambient paths, without assuming that their later
portions remain in the chart. Local finite metric lower bounds,
together with \eqref{eq:exit}, separate distinct points.
The short local paths give the reverse topological continuity.
Thus \(d_0\) induces precisely the topology of \(Z\).
Projective compactness gives metric compactness; the path definition
makes it a length metric and hence a geodesic metric.

Apply the perturbation proof of
Proposition~\ref{prop:principal-completion} to a finite subdivision
of a central actual path. Both the original and perturbed curves
are compared with the same local model, using arbitrarily small
relative errors. The simple joins have arbitrarily small total
cost. Therefore the intrinsic distance on \(U\) is \(d_0|_U\).
Density and compactness identify its completion with \(Z\).

For the last assertion choose a finite cover by connected inner
radial neighborhoods with coarse comparison at most \(2\).
Its intersection graph is connected. Choose an overlap point for
each edge in a spanning tree. Within a chart, join its center
and these points through its cone vertex, with length bounded by
its radius times the coarse comparison factor. A simple tree path
visits each chart at most once. The sum of these finitely many
local bounds, including the two endpoint joins, is a permissible
\(D_*\). No unknown global distance enters this construction.
\end{proof}

\subsection{Finite-model discrepancies}
Fix a central chart \(p\), and write \(Q_\zeta\) for its complete
finite quotient, \(d_\zeta\) for its intrinsic metric, \(r_\zeta\)
for the norm of a balanced representative, and
\(q_\zeta:Q_\zeta\to Q_0\) for the projective quotient map.
For \(r>0\), define
\begin{align}
 b_{p,r}(\theta)
 &=\sup_{\substack{q_\zeta(x)=z,\ q_\zeta(x')=z'\\
                         r_0(z),r_0(z')\le r}}
       |d_\zeta(x,x')-d_0(z,z')|,\label{eq:finite-b}\\
 a_{p,r}(\theta)
 &=\sup_{\substack{q_\zeta(x)=z\\r_0(z)\le r}}
       |r_\zeta(x)-r_0(z)|,\qquad
 \widetilde a_{p,r}(\theta)
 =\sup_{r_\zeta(x)\le r}|r_0(q_\zeta(x))-r_\zeta(x)|,
 \label{eq:finite-a}
\end{align}
where \(\zeta=\zeta_p(\theta)\); the \(d_0\) in
\eqref{eq:finite-b} is the finite zero-model distance.
We distinguish it from the central global distance by its arguments.

\begin{lemma}\label{lem:discrepancies}
The three quantities in \eqref{eq:finite-b}--\eqref{eq:finite-a}
are finite and tend to zero as \(\theta\to0\).
\end{lemma}

\begin{proof}
Properness of the projective quotient and finite invariant
properness bound balanced representatives uniformly for small
levels and bounded target radii. The source-radius condition
already gives bounded representatives in the other supremum.
All relevant sets are compact.

If one convergence failed, take a contradicting sequence and
convergent subsequences of its balanced representatives.
Source and target have the same affine invariants. At zero level
uniqueness of the balanced orbit identifies their limits.
Continuity of radius proves the two radial assertions. Finite
intrinsic parameter continuity
\cite[Theorem~5.1]{wuebben14}, applied to the two pairs of
representatives, proves the distance assertion. There are finitely
many parameter faces; passage to a face subsequence covers all
approaches.
\end{proof}

These are variational functions of specified finite moment equations,
norms and path distances. No semialgebraicity of those distances
is required. The two radial functions serve different purposes:
one bounds the preimage of a target ball, the other the image of
a source ball.

\subsection{From short pairs to global distortion}
For a correspondence \(R\subset X\times Y\), write
\[
 \dis R=\sup_{(x,y),(x',y')\in R}
          |d_X(x,x')-d_Y(y,y')|.
\]

\begin{lemma}\label{lem:short-correspondence}
Let \(X,Y\) be nonempty compact length spaces and \(R\) a
correspondence. Suppose \(\ell>0\), \(\alpha,\beta\ge0\), and
\[
 \begin{aligned}
 d_X(x,x')\le\ell&\ \Longrightarrow\
 d_Y(y,y')\le(1+\alpha)d_X(x,x')+\beta,\\
 d_Y(y,y')\le\ell&\ \Longrightarrow\
 d_X(x,x')\le(1+\alpha)d_Y(y,y')+\beta
 \end{aligned}
\]
for all related pairs. With \(D=\diam X\) and
\(q=\alpha+\beta/\ell\), one has
\begin{equation}\label{eq:short-bound}
 \dis R\le qD+\beta,\qquad
 d_{\rm GH}(X,Y)\le\tfrac12(qD+\beta),\qquad
 \diam Y\le(1+q)D+\beta .
\end{equation}
\end{lemma}

\begin{proof}
Subdivide an \(X\)-geodesic of positive length \(d_X(x,x')\)
into \(N=\lceil d_X(x,x')/\ell\rceil\) segments. Choose arbitrary
related \(Y\)-points at the interior division points and the
prescribed endpoints. Summing gives
\(d_Y(y,y')\le(1+q)d_X(x,x')+\beta\).
For zero \(X\)-distance the local inequality gives this directly.
Interchanging \(X,Y\) gives the reverse global inequality.
If \(d_Y\ge d_X\), their difference is at most \(qd_X+\beta\le qD+\beta\).
If \(d_X\ge d_Y\), it is at most \(qd_Y+\beta\le qD+\beta\).
This proves the distortion estimate, and the other assertions
follow from the correspondence bound and the first global inequality.
\end{proof}

We now specify every quantity in the global bound. Choose
\(0<\epsilon_j\downarrow0\), with \(\epsilon_j<1/8\).
At each central point choose an outer coefficient radius \(16r\)
where the central error in \eqref{eq:relative} is at most
\(\epsilon_j/2\). Extract a finite cover by central radius-\(r/4\)
images; write its radii as \(r_{ji}\).
After that choice, shrink a parameter neighborhood \(\mathcal P_j\)
so the nearby outer errors are at most \(\epsilon_j\), whole-preimage
coverage holds, and
\[
 a_{i,r_{ji}},\ \widetilde a_{i,r_{ji}}<r_{ji}/16,
 \qquad b_{i,r_{ji}}<r_{ji}.
\]
Choose any
\begin{equation}\label{eq:global-parameters}
 \begin{gathered}
 0<\ell_j<\min_i\frac{r_{ji}}{16\sqrt2},\qquad
 \alpha_j=\sqrt{\frac{1+\epsilon_j}{1-\epsilon_j}}-1,\\
 \beta_j(\theta)=\sqrt{1+\epsilon_j}\,
                  \max_i b_{i,r_{ji}}(\theta).
 \end{gathered}
\end{equation}

\begin{figure}[tb]
 \centering
 \resizebox{\linewidth}{!}{\begin{tikzpicture}[>=Latex,every node/.style={font=\small}]
 \draw[thick,->] (0,0) -- (11.2,0) node[right] {radius};
 \foreach \x/\lab in {0/0,1.6/{r/4},3/{r/2},5.1/r,10.1/{16r}} {
  \draw (\x,.1)--(\x,-.1);
  \node[below=4pt] at (\x,0) {\(\lab\)};
 }
 \draw[very thick,blue] (.15,.5) -- (1.6,.5);
 \node[above] at (1,.5) {central start};
 \draw[very thick,blue,->] (1.6,1.35) -- (3,1.35);
 \node[above,align=center] at (2.3,1.35)
    {a path of length \(\ell<r/(16\sqrt2)\)\\cannot reach \(r/2\)};
 \draw[thick,<->] (5.1,.6) -- (10.1,.6);
 \node[above,align=center] at (7.6,.6)
    {outer buffer\\first-exit cost \(>10r\)};
 \node[align=center] at (6.9,2.0)
    {whole-preimage coverage\\source start \(<r/3\)};
 \draw[dashed,->] (5.5,1.8) to[out=190,in=20] (2.1,.2);
\end{tikzpicture}
}
 \caption{The radius buffers for one chart, schematically and
 not to scale. A short pair starts in the central \(r/4\)
 ball and remains in the range where both finite discrepancies
 apply. The outer \(16r\) chart excludes ambient shortcuts.
 The source start bound uses \(a_r\); the reverse transfer
 of its endpoint uses \(\widetilde a_r\).}
 \label{fig:buffers}
\end{figure}

\begin{theorem}\label{thm:quantitative}
Under \textup{(G1)--(G3)}, let \(R_\theta\) be the graph of
\(\pi_\theta\), with the harmonic distance on \(M_\theta\) and
the completion distance on \(Z\). Then
\begin{equation}\label{eq:global-distortion}
 \dis R_\theta\le
 \inf_{j:\,\theta\in\mathcal P_j}
 \left[\left(\alpha_j+\frac{\beta_j(\theta)}{\ell_j}\right)D_*
              +\beta_j(\theta)\right].
\end{equation}
For each fixed \(j\), \(\mathcal P_j\) contains all sufficiently
small generic parameters and \(\beta_j(\theta)\to0\).
The right side tends to zero through arbitrary adjacent chambers.
Consequently \((M_\theta,d_\theta)\) converges in compact
Gromov--Hausdorff topology to the harmonic completion \((Z,d_0)\).
\end{theorem}

\begin{proof}
Fix \(j\) and one of its charts, and abbreviate \(r=r_{ji}\),
\(\epsilon=\epsilon_j\). First suppose both central endpoints have
model radius at most \(r\). Their finite central distance is at
most \(2r\), by joining through the cone vertex. Their finite
source distance is at most \(2r+b<3r\), and source radii are
at most \(r+a<17r/16\).
Finite near-minimizers stay inside the radius-\(16r\) chart,
because radius is 1-Lipschitz. Their actual images give competitors
of length less than \(4r\). On the other hand, any actual path
leaving that outer chart pays more than \(10r\) before its first
exit, by \eqref{eq:exit}. The closed sublevel images and
whole-preimage coverage justify this barrier for ambient paths.
Thus taking infima in both directions gives
\begin{equation}\label{eq:local-ambient}
 \begin{aligned}
 d_\theta(x,x')&\le(1+\alpha_j)d_0(\pi_\theta x,\pi_\theta x')
                         +\beta_j(\theta),\\
 d_0(\pi_\theta x,\pi_\theta x')
              &\le(1+\alpha_j)d_\theta(x,x')+\beta_j(\theta).
 \end{aligned}
\end{equation}
The relative factor compares the two squared-metric errors;
the larger required additive factor is \(\sqrt{1+\epsilon}\,b\).

It remains to put every short ambient pair in such a chart.
Choose an inner radius-\(r/4\) chart containing \(\pi_\theta x\).
Its entire source preimage has radius at most \(5r/16<r/3\).
If the central distance to \(\pi_\theta x'\) is at most \(\ell_j\),
a central near-minimizer cannot reach radius \(r/2\): its exit
cost is at least \(r/(4\sqrt2)>\ell_j\).
Hence the other target endpoint is inside radius \(r/2\).
If instead the source distance is at most \(\ell_j\), the source
exit cost from radius at most \(r/3\) to radius \(r/2\) is at least
\(r/(6\sqrt2)>\ell_j\). Thus \(x'\) has source radius less than
\(r/2\). The function \(\widetilde a\) then puts its target
inside radius \(r\). In either case both inequalities
\eqref{eq:local-ambient} apply.

Use Lemma~\ref{lem:short-correspondence} with
\(X=(Z,d_0)\), \(Y=(M_\theta,d_\theta)\), and the transposed
graph of \(\pi_\theta\). Both spaces are compact length spaces
and \(\diam X\le D_*\) by
Proposition~\ref{prop:central-completion}.
This proves the expression in brackets for each admissible \(j\).
Taking its infimum proves \eqref{eq:global-distortion}.

For fixed \(j\), the finite discrepancy lemma gives
\(\beta_j(\theta)\to0\). Therefore the limit superior of the
bound is at most \(\alpha_jD_*\). Now let \(j\to\infty\).
The correspondence estimate in \eqref{eq:short-bound} proves
the stated convergence.
\end{proof}

The construction first fixes accuracy, then selects a finite
cover, then takes the parameter limit. It gives an explicit
variational modulus, with no assertion of a universal power law.
In particular it does not use qualitative global
Gromov--Hausdorff convergence as a premise.
